\section*{अध्याय \ref{ch:b3:origin-of-elements} --- तत्त्व कहाँ से आते हैं}
% Values from the Book 5 scratch calculator calc/ch30.py.

\begin{solution}{exo:b3:origin-of-elements:1}
(क)~$\mathrm{Si/Fe} = 10^{7.51 - 7.46} = 10^{0.05} = 1.12$ और
$\mathrm{O/Fe} = 10^{8.69 - 7.46} = 10^{1.23} = 17.0$: लगभग एक सिलिकॉन
परमाणु और सत्रह ऑक्सीजन परमाणु प्रति लौह परमाणु। (ख)~$\mathrm{Fe/Au} =
10^{7.46 - 0.91} = 10^{6.55} = \num{3.55e6}$ और $\mathrm{Pb/Au} =
10^{1.95 - 0.91} = 10^{1.04} = 11.0$: लौह शिखर से परे प्रचुरताएँ
लाखों गुना गिरती हैं, परंतु सीसा, अपने ही शिखर पर, सोने से ग्यारह गुना
अधिक प्रचुर है।
\end{solution}

\begin{solution}{exo:b3:origin-of-elements:2}
(क)~$10^{7.46 - 12} = \num{2.884e-5}$, अर्थात् \num{28.8} लौह परमाणु प्रति
दस लाख हाइड्रोजन परमाणु। (ख)~सौ गुना कम: \num{0.288} प्रति
दस लाख, एक लौह परमाणु प्रति $1/\num{2.884e-7} = \num{3.47e6}$ हाइड्रोजन
परमाणु। (ग)~$[\mathrm{Fe/H}] = \log_{10}(\num{e-7}/\num{2.884e-5}) =
-2.46$।
\end{solution}

\begin{solution}{exo:b3:origin-of-elements:3}
$^{124}\mathrm{Xe}$ और $^{126}\mathrm{Xe}$ प्रोटॉन-समृद्ध पक्ष पर स्थित हैं,
s पथ से हटकर, और स्थायी $^{124}\mathrm{Te}$ और
$^{126}\mathrm{Te}$ द्वारा परिरक्षित हैं: \omterm{def:b3:origin-of-elements:p-process}{p नाभिक}। $^{128}\mathrm{Xe}$ और $^{130}\mathrm{Xe}$
s पथ पर हैं, परंतु
\omterm{def:b3:origin-of-elements:s-process}{r प्रक्रम} से स्थायी $^{128}\mathrm{Te}$ और $^{130}\mathrm{Te}$ द्वारा परिरक्षित: केवल s। $^{129}\mathrm{Xe}$,
$^{131}\mathrm{Xe}$ और $^{132}\mathrm{Xe}$ s पथ पर हैं और इनका
निम्न $Z$ वाला कोई स्थायी समभारिक नहीं है (आयोडीन में केवल 127 है): दोनों प्रक्रम।
$^{134}\mathrm{Xe}$ और $^{136}\mathrm{Xe}$ उस सीमा के पार स्थित हैं,
$^{133}\mathrm{Xe}$ के आगे, जहाँ s पथ ज़ेनॉन से हट जाता है, और
परिरक्षित नहीं हैं: केवल r।
\end{solution}

\begin{solution}{exo:b3:origin-of-elements:4}
(क)~$t = d/c = 40\times\qty{3.086e22}{m}/\qty{2.998e8}{m.s^{-1}} =
\qty{4.12e15}{s}$, लगभग \num{130} मिलियन वर्ष। (ख)~$\Delta v/c \le
1.7/\num{4.12e15} = \num{4.1e-16}$। यह एक ऊपरी सीमा है, क्योंकि विलंब का
कुछ या सारा भाग गामा किरणों के उत्सर्जन से आ सकता है, जो
विलय के बाद होता है।
\end{solution}

\begin{solution}{exo:b3:origin-of-elements:5}
(क)~$295\times0.1067 = \qty{31.48}{mb}$, $832\times0.0373 =
\qty{31.03}{mb}$, $155\times0.1984 = \qty{30.75}{mb}$। (ख)~
$^{122}\mathrm{Te}$ से: $N_{123} = 31.48/832 = \num{0.03783}$, \qty{1.4}{\percent}
मापे गए मान से ऊपर, और $N_{124} = 31.48/155 = \num{0.2031}$,
\qty{2.4}{\percent} ऊपर। जो प्रचुरताएँ पाँच गुने के गुणक से भिन्न हैं,
उनका अनुमान केवल क्रॉस-सेक्शन से ही कुछ प्रतिशत के भीतर लगाया जा सकता है।
(ग)~$^{128}\mathrm{Te}$ केवल r है: s प्रवाह टेल्यूरियम को
$^{127}\mathrm{Te}$ पर छोड़ देता है, जो क्षय होता है, इसलिए कोई स्थायी s प्रवाह
$^{128}\mathrm{Te}$ से होकर नहीं गुजरता, और इसकी प्रचुरता \omterm{def:b3:origin-of-elements:s-process}{r प्रक्रम} से निर्धारित हुई।
\end{solution}

\begin{solution}{exo:b3:origin-of-elements:6}
(क)~$\tau_0\langle\sigma\rangle = 0.306\times155 = 47.4$, अतः
$f = 1/(1 + 1/47.4) = \num{0.979}$, $^{124}\mathrm{Te}$ के लिए;
$^{138}\mathrm{Ba}$ के लिए, $0.306\times4.0 = 1.224$ और $f = 1/(1 + 0.817) =
\num{0.550}$: जादुई नाभिक केवल \qty{55}{\percent} तक का
प्रवाह आगे बढ़ाता है। (ख)~छह गुणक \num{0.959}, \num{0.550}, \num{0.908},
\num{0.771}, \num{0.972} और \num{0.915} हैं; इनका गुणनफल \num{0.328} है:
$\sigma N$ \num{3.05} गुना घटता है। (ग)~मापी गई कमी
$27.6/7.72 = 3.57$ है। मॉडल चरण के आकार को
\qty{15}{\percent} के भीतर ठीक बैठाता है, जिसका लगभग सारा भाग $^{138}\mathrm{Ba}$ और
$^{140}\mathrm{Ce}$ से आता है, ये दो नाभिक हैं जिनका $N = 82$ है।
\end{solution}

\begin{solution}{exo:b3:origin-of-elements:7}
(क)~$kT = 8\times10^3\times\qty{1.602e-19}{J} = \qty{1.282e-15}{J}$, और
$v_T = \sqrt{2\times\num{1.282e-15}/\num{1.6749e-27}} =
\qty{1.24e6}{m.s^{-1}}$। (ख)~$\langle\sigma\rangle =
\qty{155e-31}{m^2}$ के साथ, $\lambda = n_\mathrm{n}\langle\sigma\rangle v_T$
\qty{1.92e-11}{s^{-1}} है, $n_\mathrm{n} = \qty{e12}{m^{-3}}$ के लिए
($10^6$ प्रति घन सेंटीमीटर): माध्य समय \qty{5.2e10}{s}, लगभग
\num{1650} वर्ष; फिर $10^8$ पर \num{16.5} वर्ष और \num{0.165}~वर्ष,
दो महीने, \qty{e10}{cm^{-3}} पर। (ग)~$^{127}\mathrm{Te}$ घंटों में
क्षय होता है, इन सभी ग्रहण कालों से कहीं तेज़: यह सदैव क्षय होता है और
पथ वहाँ नहीं बँटता। $^{134}\mathrm{Cs}$ लगभग दो
वर्ष जीवित रहता है, $10^8$ और \qty{e10}{cm^{-3}} पर ग्रहण कालों के बीच:
हीलियम फ़्लैश के उच्च घनत्वों पर इसका अच्छा-खासा भाग
क्षय से पहले ग्रहण हो जाता है। यह एक शाखा-बिंदु है, और इसके दोनों ओर की प्रचुरताएँ
न्यूट्रॉन घनत्व को मापती हैं।
\end{solution}

\begin{solution}{exo:b3:origin-of-elements:8}
(क)~कैल्शियम से $10^{6.30 - 3.14} = 1445$ गुना दुर्लभ और
$10^{4.97 - 3.14} = 68$ गुना दुर्लभ टाइटेनियम से। (ख)~$10^{8.46 - 1.38}
= \num{1.2e7}$। (ग)~स्कैंडियम विषम--सम प्रभाव दिखाता है: प्रोटॉनों की सम संख्या
वाले नाभिक अधिक दृढ़ता से बँधे होते हैं, और उनका निर्माण विशाल तारों के दहन चरणों में
हीलियम नाभिकों के ग्रहण से होता है, जो
इन्हें अनुकूल होते हैं। बेरिलियम में लिथियम, बेरिलियम और बोरॉन की गर्त दिखती है:
ये नाजुक नाभिक तारकीय अंतरभागों के तापमानों पर नष्ट होते हैं, और
हीलियम दहन सीधे हीलियम से कार्बन तक जाता है,
इन्हें बनाए बिना।
\end{solution}

\begin{solution}{exo:b3:origin-of-elements:9}
(क)~मान लीजिए $\alpha$ \omterm{def:b3:origin-of-elements:alpha}{अल्फ़ा तत्त्व} हैं और $\mathrm{Fe}_\mathrm{cc}$,
$\mathrm{Fe}_\mathrm{Ia}$ दो प्रकार के अधिनवतारों से आया लौह है। पठार पर
$\alpha/\mathrm{Fe}_\mathrm{cc} = 10^{P}$ गुना सौर
अनुपात। $[\mathrm{Fe/H}] = 0$ पर, $\alpha/(\mathrm{Fe}_\mathrm{cc} +
\mathrm{Fe}_\mathrm{Ia})$ सौर है। भाग देने पर,
$\mathrm{Fe}_\mathrm{cc}/(\mathrm{Fe}_\mathrm{cc} + \mathrm{Fe}_\mathrm{Ia})
= 10^{-P}$, और टाइप~Ia का भाग $1 - 10^{-P} = 1 - 10^{-0.4} =
\num{0.602}$। (ख)~$10^{-P} = 1/3$ से प्राप्त होता है $P = \log_{10}3 = 0.48$।
(ग)~टाइप~Ia अधिनवतारे दोनों मंदाकिनियों में एक ही विलंब के बाद आरंभ होते हैं।
जब ये वामन मंदाकिनी में आरंभ हुए, तब उसकी गैस केवल
$[\mathrm{Fe/H}] = -2$ तक पहुँची थी, जबकि आकाशगंगा में $-1$ था: उसने अपनी गैस को
तारों में दस गुना धीमी गति से बदला था।
\end{solution}

\begin{solution}{exo:b3:origin-of-elements:10}
(क)~ज़िरकोनियम के लिए,
\[
\frac{(X/\mathrm{Mg})_\text{पृथ्वी}}{(X/\mathrm{Mg})_\mathrm{CI}} =
\frac{10.5/228000}{3.96/95870} = 1.11 ;
\]
इसी प्रकार सोडियम के लिए $(2670/228000)/(5010/95870) = 0.224$,
पोटैशियम के लिए \num{0.190} और रूबिडियम के लिए \num{0.118}। (ख)~सबसे अधिक दुर्गलनीय ज़िरकोनियम
अपने कॉन्ड्राइट अनुपात के निकट है। तीन मध्यम वाष्पशील
क्षारीय धातुएँ, जो \qtyrange{800}{1006}{K} के बीच संघनित होती हैं, पाँच से नौ गुणक
तक घटी हुई हैं; रूबिडियम, जो सबसे अधिक वाष्पशील है, सबसे अधिक घटा है। सिलिकेट पृथ्वी
\omterm{def:b3:origin-of-elements:volatility}{वाष्पशीलता प्रवृत्ति} का अनुसरण करती है। (ग)~मैग्नीशियम एक प्रमुख,
दुर्गलनीय तत्त्व है जो चट्टान के अधिकांश भाग के साथ संघनित होता है, और यह
लिथोफ़िल है: क्रोड बनते समय यह प्रावार में ही रह गया। एक सिडरोफ़िल
संदर्भ, जैसे लौह या निकल, क्रोड की ओर हुई हानि को
वाष्पशीलता से हुई हानि के साथ मिला देगा।
\end{solution}

\begin{solution}{exo:b3:origin-of-elements:11}
(क)~जब द्रव्यमान $\mathrm{d}s$ तारों में बंद हो जाता है, तब पवन
गैस का $\lambda\,\mathrm{d}s$ हटा देती है, जिसकी धातुता $Z$ है:
$\mathrm{d}M_\mathrm{g} = -(1 + \lambda)\,\mathrm{d}s$ और
$\mathrm{d}(ZM_\mathrm{g}) = y\,\mathrm{d}s - Z\,\mathrm{d}s -
\lambda Z\,\mathrm{d}s$। दूसरे व्यंजक को विस्तारित करके
$Z\,\mathrm{d}M_\mathrm{g}$ घटाने पर $M_\mathrm{g}\,\mathrm{d}Z =
y\,\mathrm{d}s = -\frac{y}{1 + \lambda}\,\mathrm{d}M_\mathrm{g}$ प्राप्त होता है,
जिसका समाकलन परिणाम देता है। (ख)~$\ln(1/0.29) = 1.238$: $Z = 0.0155$
$\lambda = 0$ के लिए, और उसका आधा, $0.0077$, $\lambda = 1$ के लिए। (ग)~पवन
प्रभावी लब्धि को गुणक $1 + \lambda$ से घटा देती है। कम द्रव्यमान की
मंदाकिनी में अधिनवतारे गैस को अधिक आसानी से बाहर धकेलते हैं, $\lambda$ बड़ा होता है,
और धातुएँ गैस के साथ निकल जाती हैं।
\end{solution}

\begin{solution}{exo:b3:origin-of-elements:12}
(क)~सूर्य में $N_\mathrm{Li}/N_\mathrm{H} = 10^{0.96 - 12} = \num{9.1e-12}$
और CI कॉन्ड्राइट में $10^{3.25 - 12} = \num{1.78e-9}$ है। सूर्य
अपने लिथियम का $10^{-2.29} = \num{0.0051}$ भाग बचाए रखता है: उसने
उसका \qty{99.5}{\percent} नष्ट कर दिया है। (ख)~$\num{1.78e-9}/\num{1.45e-10} = 12.3$ और
$\num{1.78e-9}/\num{4.72e-10} = 3.8$। (ग)~सौरमंडल में बिग बैंग द्वारा बने लिथियम से लगभग चार गुना अधिक लिथियम है:
उसका अधिकांश लिथियम
बाद में किसी अन्य प्रक्रम से बना। सबसे पुराने तारों में बिग बैंग द्वारा बने
लिथियम का तीन गुना कम है: या तो उन्होंने उसका कुछ भाग नष्ट कर दिया है, जैसा
सूर्य ने किया है, या पूर्वानुमान में कुछ छूट गया है; केवल ये दो तुलनाएँ
अकेले यह नहीं बतातीं कि इनमें से कौन सा है।
\end{solution}

\begin{solution}{pb:b3:origin-of-elements:1}

\textbf{भाग I.}

\textbf{1.} $10^{12 - 7.46} = 10^{4.54} = \num{34700}$।

\textbf{2.} $N_\mathrm{He}/N_\mathrm{H} = 10^{10.914 - 12} = \num{0.0820}$, इसलिए
$4N_\mathrm{He}/N_\mathrm{H} = \num{0.328}$, जबकि $Y/X =
0.2423/0.7438 = \num{0.326}$। ये \qty{0.7}{\percent} तक मेल खाते हैं; छोटा
अंतर इसलिए है क्योंकि एक हीलियम परमाणु
चार हाइड्रोजन परमाणुओं से थोड़ा कम भारी होता है।

\textbf{3.} $\log_{10}\bigl((1/350000)/\num{2.884e-5}\bigr) = -5.544 + 4.540 =
-1.00$।

\textbf{4.} $[\mathrm{O/H}] = -1.00 + 0.40 = -0.60$, इसलिए
$\log\varepsilon(\mathrm{O}) = 8.69 - 0.60 = 8.09$। तारे में प्रति लौह परमाणु
$10^{8.09 - 6.46} = 42.7$ ऑक्सीजन परमाणु हैं, जबकि
सूर्य में $10^{1.23} = 17.0$ हैं: यानी $10^{0.4} = 2.5$ गुना अधिक।

\textbf{5.} $Z \approx 0.0139\times10^{-0.60} = \num{0.0035}$, सूर्य के
मान का एक चौथाई।

\textbf{6.} ठंडे तारे के स्पेक्ट्रम में लौह की बहुत अधिक रेखाएँ होती हैं, जबकि
ऑक्सीजन की रेखाएँ कम और मंद होती हैं: लौह कहीं अधिक आसानी से
और सूक्ष्मता से मापा जाता है, और धातुओं के प्रतिनिधि के रूप में काम आता है।

\textbf{भाग II.}

\textbf{7.} प्रत्येक कम $Z$ वाले स्थिर समदाब रेखा के पीछे स्थित है:
$^{122}\mathrm{Sn}$, $^{123}\mathrm{Sb}$ और $^{124}\mathrm{Sn}$ \omterm{def:b3:origin-of-elements:s-process}{r प्रक्रम} की
क्षय श्रृंखलाओं को टेल्यूरियम तक पहुँचने से पहले रोक देते हैं। ये तीनों उस s पथ पर स्थित हैं,
जो टेल्यूरियम से होकर $^{122}\mathrm{Te}$ से
$^{126}\mathrm{Te}$ तक जाता है।

\textbf{8.} \num{31.48}, \num{31.03} और \qty{30.75}{mb}; माध्य
\qty{31.09}{mb}।

\textbf{9.} अन्य दो का माध्य $(31.48 + 30.75)/2 = \qty{31.11}{mb}$ है,
इसलिए $N_{123} = 31.11/832 = \num{0.03740}$, जो मापे गए \num{0.0373} से
केवल \qty{0.3}{\percent} ऊपर है।

\textbf{10.} $1.329\times44.4 = \qty{59.0}{mb}$, जो पठार का लगभग दोगुना है।
$^{128}\mathrm{Te}$ केवल \omterm{def:b3:origin-of-elements:s-process}{r प्रक्रम} से बना है: s पथ टेल्यूरियम को
$^{127}\mathrm{Te}$ पर छोड़ता है, इसलिए इसकी प्रचुरता \omterm{def:b3:origin-of-elements:s-process}{r प्रक्रम} से बनी है, और
उसमें s प्रवाह का कोई योगदान नहीं।

\textbf{11.} $1 - 7.43/8.53 = \num{0.13}$: प्रवाह का \qty{13}{\percent} भाग
$^{148}\mathrm{Sm}$ के चारों ओर से होकर गया, अस्थिर
$^{147}\mathrm{Nd}$, $^{147}\mathrm{Pm}$ और $^{148}\mathrm{Pm}$ पर ग्रहण द्वारा, इससे पहले कि
वे क्षय होकर समैरियम बन पाते।

\textbf{12.} 82 न्यूट्रॉनों के बंद कोश के कारण ये नाभिक एक अतिरिक्त
न्यूट्रॉन को कमज़ोर रूप से बाँधते हैं और इनके \omterm{def:b3:origin-of-elements:s-process}{न्यूट्रॉन प्रग्रहण} अनुप्रस्थ-परिच्छेद बहुत छोटे हैं (\qty{4.0}{mb}
$^{138}\mathrm{Ba}$ के लिए)। प्रवाह इन पर जमा होता है और उसका केवल एक भाग
आगे निकलता है: ये अवरोध-बिंदु हैं।

\textbf{भाग III.}

\textbf{13.} $\mu_1 = 13.7/(13.7 + 33.4) = 13.7/47.1 = \num{0.291}$।

\textbf{14.} $\ln(1/0.291) = \num{1.235}$, इसलिए $y = 0.0154/1.235 = \num{0.0125}$,
$Z_\odot$ का \num{0.81} गुना।

\textbf{15.} $\mu = 0.9$ पर $\ln(1/0.9) = 0.1054$ है, इसलिए $Z = 0.0125\times0.1054 =
\num{0.00131}$। चूँकि $0.00131/0.0154 = 0.0851$ है, $[\mathrm{Fe/H}] =
-1.07$। $\mu = 0.5$ पर $Z = 0.0125\times0.693 = \num{0.00864}$ और
$[\mathrm{Fe/H}] = -0.25$।

\textbf{16.} $Z = 0.00154$ है, इसलिए $\mu = \mathrm{e}^{-0.00154/0.01247} =
\mathrm{e}^{-0.1235} = \num{0.884}$।

\textbf{17.} $1 - 0.884 = \num{0.116}$: द्रव्यमान का \qty{12}{\percent}।

\textbf{18.} $y = 0.0154/\ln 10 = \num{0.0067}$, लगभग आधा।

\textbf{भाग IV.}

\textbf{19.} $Z$ से कम \omterm{def:b3:origin-of-elements:metallicity}{धात्विकता} वाले तारों का द्रव्यमान
$M(1 - \mathrm{e}^{-Z/y})$ है, इसलिए प्रति एकक $Z$ यह
$(M/y)\,\mathrm{e}^{-Z/y}$ है। चूँकि $Z = Z_\odot10^{[\mathrm{Fe/H}]}$ है,
$\mathrm{d}Z = Z\ln 10\,\mathrm{d}[\mathrm{Fe/H}]$, और प्रति एकक
$[\mathrm{Fe/H}]$ के प्रति एकक बंटन समानुपाती है।
$Z\mathrm{e}^{-Z/y}$ है। इसका अवकलज,
$(1 - Z/y)\mathrm{e}^{-Z/y}$, $Z = y$ पर शून्य हो जाता है: शिखर
$[\mathrm{Fe/H}] = \log_{10}0.810 = -0.09$ पर है।

\textbf{20.} $\num{0.116}/(1 - 0.291) = 0.1162/0.7091 = \num{0.164}$।

\textbf{21.} $Z/y = 1.235\times10^{-0.5} = 0.3905$, $1 - \mathrm{e}^{-0.3905} =
0.3233$, और $0.3233/0.7091 = \num{0.456}$।

\textbf{22.} $6670\times0.1638 = 1093$, जबकि 26 है: लगभग चालीस गुना अधिक।

\textbf{23.} $26/902 = \qty{2.9}{\percent}$, फिर भी बॉक्स के अनुमान से
पाँच गुना से अधिक कम है। मॉडल के तारे दीर्घजीवी हैं, वे जिनका द्रव्यमान
अधिकतम लगभग एक सौर द्रव्यमान है, और जो बनने के बाद से सब बचे हुए हैं;
अधिक द्रव्यमान वाले F तारों में सबसे पुराने, जो सबसे कम धात्विक भी हैं,
मर चुके हैं, और पूरा नमूना उन्हें कम गिनता है।

\textbf{24.} $(1 - 0.1^{0.1})/(1 - 0.1) = 0.2057/0.9 = \num{0.23}$: और खराब।
गैस-भिन्नता बढ़ाने से भी मदद नहीं मिलती: जैसे-जैसे $\mu_1$ 1 की ओर जाता है,
भिन्नता $0.1$ की ओर जाती है, जो अब भी देखे गए \qty{2.9}{\percent} से तीन गुना
से अधिक है। कोई भी गैस-भिन्नता बंद बॉक्स को नहीं बचाती।

\textbf{25.} बंद बॉक्स यह भविष्यवाणी करता है कि सौर पड़ोस के
\textbf{\qty[text-series-to-math=true]{16}{\percent}} दीर्घजीवी तारों का
$[\mathrm{Fe/H}] < -1$ होगा, जबकि देखा गया अधिकतम \qty{3}{\percent} है।
चक्रिका बंद बॉक्स नहीं थी: वह कम धात्विक गैस के धीमे अंतर्वाह से
बढ़ी, इसलिए निम्न \omterm{def:b3:origin-of-elements:metallicity}{धात्विकता} पर तारे बनाने के लिए बहुत कम गैस कभी
प्रतीक्षा में बैठी।
\end{solution}
